\index{Cartesian coordinate system}\index{Cartesian coordinates}
To locate the position of the \robot\ on the screen, we use a {\em Cartesian coordinate system}
as shown in figure \ref{fig:repcar} below:

\begin{figure}[htbp]
 \includevspacePostScript[6.2cm]{RepereCartesien}{0.95}{0.95}
 \caption{Coordinates in a Cartesian coordinate system}\label{fig:repcar}
\end{figure}

The {\em origin\/} $0$ of the coordinate system is located at the centre of the grid. The horizontal axis $OX$ is 
called the {\em x-axis\/} (axis of {\em abscissas\/})\index{Abscissa}. The vertical axis $OY$ is 
called the {\em y-axis\/} (axis of {\em ordinates\/})\index{Ordinate}. Both axes are graduated. The length of the
graduations on the x-axis is equal to the length of the horizontal side of a square 
 of the grid counted in pixels. Likewise, the length of the
graduations on the y-axis is equal to the length of the vertical side of a square
of the grid counted in pixels. 

A grid point is located by its {\em  Cartesian coordinates}\index{Cartesian coordinates}, which are
two integers $(X, Y)$, the first being its {\em abscissa} (or x-coordinate) and the second its {\em ordinate} (or y-coordinate). For example the 
grid point located at the centre of the grid has coordinates $(0, 0)$. The points located
to the right of the origin $O$ have a positive abscissa. The points located
to the left of the origin $O$ have a negative abscissa (preceded by the minus sign $-$). 
Likewise, the points located above
the origin $O$ have a positive ordinate while the points located
below the origin $O$ have a negative ordinate.  As an example, some points of 
the grid are labelled with their coordinates in figure \ref{fig:repcar}  (page \pageref{fig:repcar}).  

The abscissa X of a point A is the distance from the origin of its projection on the axis $OX$
(that is, of the point of intersection
of the axis $OX$ with the line parallel to the axis $OY$ passing through A). Likewise, the ordinate Y of a point 
A is the distance from the origin of its projection on the axis $OY$
(that is, of the point of intersection
of the axis $OY$ with the line parallel to the axis $OX$ passing through A). For example in figure
\ref{fig:coorA} 
below,
we have $X = 6$ and $Y = 7$\footnote{\ In the folder {\idt book/\allowbreak 14\_cartesian\_coordinates} you will find
a \python\ program, {\idt Learn\_Coordinates.py}, for learning how to locate points of the plane with
a Cartesian coordinate system.}.

\begin{figure}[htbp]
 \includevspacePostScript[4.1cm]{AbscisseOrdonnee}{0.95}{0.95}
 \caption{Point A with abscissa X = 6 and ordinate Y = 7}
 \label{fig:coorA}
\end{figure}

The command {\idt moveTo(x,\ y)} places the \robot\ on the grid point with Cartesian coordinates $(x, y)$. 
The orientation of the \robot\ and the position of its pen remain unchanged. 

For example, the command {\idt moveTo(0,\ 0)} places the \robot\ at the origin of the Cartesian
coordinate system, at the centre of the grid\footnote{\ The command {\idt gc()} has the additional effect
of recentring, if necessary, the origin of the coordinate system exactly at the centre of the screen.}. 
An example is given below:

\clearpage
\vspace*{-1cm}
 {\raggedright 
 \vbox to 0pt{\vspace*{\baselineskip}\vspace*{6cm}%
 \includevspacePostScript[-\baselineskip]{LigneDeCarres}{0.95}{0.95}}%
 \label{prog:LigneDeCarres1}
 \input python/RowOfSquares1.tex
 }

The command {\idt translate(x,\ y)} translates the \robot\ by $(x, y)$
without changing its orientation\footnote{\ In the folder {\idt book/\allowbreak 14\_cartesian\_coordinates} you will find
a \python\ program, {\idt Learn\_Translation.py}, for learning how to translate points of the plane 
located by Cartesian coordinates.}. An example of the translation of a rectangle is given in figure 
\ref{fig:extrans} below.
\begin{figure}[h]
\includevspacePostScript[6.0cm]{Translater}{0.95}{0.95}%
 \caption{Translation (17, -8) from A to B}\label{fig:extrans}
\end{figure}
\newpage
A translation of the \robot\ by $(x, y)$ consists in increasing (or decreasing if $x$ is negative)
its abscissa by $x$ and increasing (or decreasing if $y$ is negative)
its ordinate by $y$. If the centre of the grid, with coordinates $(0, 0)$, is exactly at the centre
of the screen, we can define moving in terms of translation as follows:

{
\parindent \parindentprog
\medskip

\vspace{\neglskip}\mbox{\blc {\mcl def}{\idt \ moveTo(x,\ y):}}\par
\vspace{\neglskip}\mbox{\blc \blk\blk{\idt gc();\ translate(x,\ y)}}\par

\medskip }

\noindent As this example shows, the commands {\idt moveTo} and {\idt translate}
are procedures with two parameters, the abscissa {\idt x} then the ordinate {\idt y},
placed between parentheses and separated by a comma, like the two parameters
of the call {\idt rectangle(3,\ 4)} in the previous chapter. The order matters:
{\idt moveTo(6,\ 7)} and {\idt moveTo(7,\ 6)} do not place the \robot\ on the same point.


The program {\idt RowOfSquares1} on page
\pageref{prog:LigneDeCarres1} can be rewritten us\-ing a
translation command instead of a move command,
as follows:\\
{\raggedright%
\vspace{\baselineskip}%
\vbox to 0pt{\vspace*{\baselineskip}\vspace*{6cm}%
\includevspacePostScript[-\baselineskip]{LigneDeCarres}{0.95}{0.95}}%
\vspace*{-1cm}%
\input python/RowOfSquares2.tex
\vspace{0.5cm}% the figure placed beside goes further down than the Python listing
}
 
 With the command {\idt translate(x,\ y)}, the movement of the \robot\ is relative to its starting
position, whereas with the command {\idt moveTo(x,\ y)}, the movement of the \robot\ is relative 
to the origin of the Cartesian coordinate system.
 
\needspace{6\baselineskip}\noindent \rule{137mm}{0.1mm}\par\nopagebreak %---------------------------------------------
\begin{exercices}{Write \python\ programs using the commands
{\idt move\-To} and {\idt trans\-late} to produce the following drawings:

\tiret  the tilings obtained by repeating various patterns inscribed in a square of
size $(6 \times 6)$ in figure \ref{fig:pm6x6} (page \pageref{fig:pm6x6});

\tiret  the tiling of quadrilaterals in figure \ref{fig:pavqua} (page \pageref{fig:pavqua});

\tiret  the mosaic of the Alhambra recorded by Escher \cite{escher1} in figure 
\ref{fig:mosa5} (page \pageref{fig:mosa5}).

\noindent $\Box$}\end{exercices}
\begin{figure}[htbp]
 \includevspacePostScript[7.7cm]{PavageQuadrilateres}{0.95}{0.95}%
 \caption{Tiling of quadrilaterals}\label{fig:pavqua}
\end{figure}
\begin{figure}[htbp]
 \includevspacePostScript[5.8cm]{MosaiqueAlhambra5}{0.95}{0.95}
 \caption{Mosaic of the Alhambra}\label{fig:mosa5}
\end{figure}
\begin{figure}[htbp]
 \includevspacePostScript[17.2cm]{PavageMotifs6x6}{0.95}{0.95}
 \caption{Tilings of $(6 \times 6)$ patterns}\label{fig:pm6x6}
\end{figure}
\begin{corrige}{
We give the program drawing the tiling of quadrilaterals of figure \ref{fig:pavqua} (page 
\pageref{fig:pavqua}), the other programs being in the folder {\idt book/\allowbreak 14\_cartesian\_coordinates}. It is a tiling of
the screen with patterns inscribed in $(6 \times 6)$ squares, the basic pattern being
drawn next to the procedure that draws it.
%\noindent\includePostScript{MotifPavageQuadrilateres}{0.95}{0.95}
\input python/QuadrilateralTiling.tex
\noindent $\Box$}\end{corrige}
\vfill

%%%%%%%%%%%%%%%%%%%%%%%%%%%%%%%%%%%%%%%%%%%%%%%%%%%%%%%%%%%%%%%%%%%%%%%%%%%
\noindent \framebox[\textwidth]{
\begin{minipage}{\myparboxwidth}
\setlength{\parindent}{2em}
\small \sf
\vspace{\abovedisplayskip}
\section*{Addenda}
 
 Following in the footsteps of Apollonius\footnote{\ \href{https://en.wikipedia.org/wiki/Apollonius_of_Perga}{Apollonius of Perga}, Greek astronomer and
 mathematician, 295 -- 230 BC.}, who had informally given the equations
 of the conics (circles, ellipses, parabolas and hyperbolas), Fermat and Descartes were the first 
 to ``put into equations'' geometrical problems. 
 Descartes introduced the notion
 of a system of plane coordinates, each point of the plane being defined by its abscissa $x$
 and its ordinate $y$, obtained by projecting onto axes $OX$ and $OY$, rectangular or oblique. 
 
 An algebraic curve is then defined by the relation $f(x, y) = 0$ that holds between the coordinates of
 each of its points. 
For example the Cartesian equation of the circle with centre $(a, b)$ is
 $x^{2} + y^{2} - 2 a x - 2 b y + c = 0$ with the radius $R$ such that $R^{2} = a^{2} + b^{2} -c$,
 assuming that $a^{2} + b^{2} \geq c$. 
 
\vspace{\abovedisplayskip}
\end{minipage}
}
%%%%%%%%%%%%%%%%%%%%%%%%%%%%%%%%%%%%%%%%%%%%%%%%%%%%%%%%%%%%%%%%%%%%%%%%%%%% 

%%%%%%%%%%%%%%%%%%%%%%%%%%%%%%%%%%%%%%%%%%%%%%%%%%%%%%%%%%%%%%%%%%%%%%%%%%%
\noindent \framebox[\textwidth]{
\begin{minipage}{\myparboxwidth}
\setlength{\parindent}{2em}
\small \sf
\vspace{\abovedisplayskip}
\indent  Descartes set out his method in a work
 published in 1637
 entitled ``La Géométrie'' (Geometry) which, together with ``La Dioptrique'' (Optics) and ``Les Météores'' (Meteorology), formed the
 appendices of one of the most famous texts of philosophy, the celebrated ``Discours de la 
 méthode pour bien conduire sa raison et chercher la vérité 
 dans les sciences'' (Discourse on the Method of Rightly Conducting One's Reason and of Seeking Truth in the Sciences), where he sets out his methodical and rational philosophy, since called
{\em  Cartesianism}. The idea of an equation appears more clearly in Fermat, who discovered
 independently of Descartes, around 1629, the basic principle of {\em analytic
 geometry}. He only published it in 1679 in his ``Isagoge (Introduction aux lieux plans et solides)'' (Introduction to plane and solid loci).
 
 Descartes had only seriously studied the conics of the second degree. 
 Newton\footnote{\ \href{https://en.wikipedia.org/wiki/Isaac_Newton}{Isaac Newton}, English mathematician, physicist and astronomer, 1642--1727.}
 studied the cubics of the third degree and curves of higher degree. We also
 owe him the systematic use of negative coordinates, of which Descartes was
 wary. La Hire\footnote{\ \href{https://en.wikipedia.org/wiki/Philippe_de_La_Hire}{Philippe de la Hire}, French astronomer and mathematician, 1640--1718}
 introduced in 1679 the consideration of a system of spatial coordinates, which
 amounts to building a  third coordinate axis $OZ$, perpendicular to the axes $OX$
 and $OY$ defining the plane. A point in space is then determined by three real
 numbers $(x, y, z)$ and the equation $f(x, y, z) = 0$ defines a skew surface. 
 Möbius\footnote{\ \href{https://en.wikipedia.org/wiki/August_Ferdinand_Möbius}{August Ferdinand Möbius}, German astronomer and mathematician,
 1790--1868} and Plücker%
 \footnote{\ \href{https://en.wikipedia.org/wiki/Julius_Plücker}{Julius Plücker}, German mathematician and physicist, 1801--1868} extended
 the notion of coordinates by introducing the homogeneous coordinates \((X, Y, Z, T)\) such that
 $x = \frac{X}{T}$,  $y = \frac{Y}{T}$ and  $z = \frac{Z}{T}$. This makes it possible to introduce the points
 at infinity $(X, Y, Z, 0)$ and the imaginary points $X = A_{1} + i A_{2}$, $Y = B_{1} + i B_{2}$,
 $Z = C_{1} + i C_{2}$ and $T = D_{1} + i D_{2}$ where $i^{2} = -1$. One can then consider for
 example imaginary circles such that $a^{2} + b^{2} < c$.
 
 \vspace{\abovedisplayskip}
\end{minipage}
}
%%%%%%%%%%%%%%%%%%%%%%%%%%%%%%%%%%%%%%%%%%%%%%%%%%%%%%%%%%%%%%%%%%%%%%%%%%%% 
\newpage                               % If the chapter ends with a blank page
\thispagestyle{empty}           % If the chapter ends with a blank page
